\documentclass[bachelor, och, pract]{../SCWorks}
% параметр - тип обучения - одно из значений:
%    spec     - специальность
%    bachelor - бакалавриат (по умолчанию)
%    master   - магистратура
% параметр - форма обучения - одно из значений:
%    och   - очное (по умолчанию)
%    zaoch - заочное
% параметр - тип работы - одно из значений:
%    referat    - реферат
%    coursework - курсовая работа (по умолчанию)
%    diploma    - дипломная работа
%    pract      - отчет по практике
%    pract      - отчет о научно-исследовательской работе
%    autoref    - автореферат выпускной работы
%    assignment - задание на выпускную квалификационную работу
%    review     - отзыв руководителя
%    critique   - рецензия на выпускную работу
% параметр - включение шрифта
%    times    - включение шрифта Times New Roman (если установлен)
%               по умолчанию выключен
\usepackage[T2A]{fontenc}
\usepackage[utf8]{inputenc}
\usepackage{graphicx}
\usepackage[hyphens]{url} 

\usepackage[sort,compress]{cite}
\usepackage{amsmath}
\usepackage{amssymb}
\usepackage{amsthm}
\usepackage{fancyvrb}
\usepackage{longtable}
\usepackage{listings}
\usepackage{array}
\usepackage[english,russian]{babel}
\usepackage{minted}
\usepackage{tempora}

\usepackage[colorlinks=true]{hyperref}

\lstset{
    extendedchars=\true,
    inputencoding=utf8x,
    language=Python,
    numbers=left,
    backgroundcolor=\color{white},
    basicstyle=\ttfamily\footnotesize,
    breaklines=true,
    frame=none,
    tabsize=4,
    xleftmargin=0pt,
    xrightmargin=0pt,
    aboveskip=0pt,
    belowskip=0pt
}

\newcommand{\eqdef}{\stackrel {\rm def}{=}}

\newtheorem{lem}{Лемма}
\worktype{Отчет о практике}
\begin{document}


% Кафедра (в родительном падеже)
\chair{математической кибернетики и компьютерных наук}

% Тема работы
\title{Проект. Приложение <<Team Finder>> }

% Курс
\course{2}

% Группа
\group{251}

% Факультет (в родительном падеже) (по умолчанию "факультета КНиИТ")
%\department{факультета КНиИТ}

% Специальность/направление код - наименование
\napravlenie{09.03.04 "--- Программная инженерия}
%\napravlenie{02.03.01 "--- Математическое обеспечение и администрирование информационных систем}
%\napravlenie{09.03.01 "--- Информатика и вычислительная техника}
%\napravlenie{09.03.04 "--- Программная инженерия}
%\napravlenie{10.05.01 "--- Компьютерная безопасность}

% Для студентки. Для работы студента следующая команда не нужна.
%\studenttitle{Студентки}

% Фамилия, имя, отчество в родительном падеже
\author{Архипова Ивана Сергеевича}

% Заведующий кафедрой
\chtitle{доцент, к.\,ф.-м.\,н.} % степень, звание
\chname{С.\,В.\,Миронов}

%Научный руководитель (для реферата преподаватель проверяющий работу)
\satitle{доцент, к.\,ф.-м.\,н.} %должность, степень, звание
\saname{И.\,А.\,Батраева}

% % Руководитель практики от организации (только для практики,
% % для остальных типов работ не используется)
% \patitle{к.\,ф.-м.\,н., доцент}
% \paname{Д.\,Ю.\,Петров}

% Семестр (только для практики, для остальных
% типов работ не используется)
\term{4}

% Наименование практики (только для практики, для остальных
% типов работ не используется)
\practtype{проектная}

% Продолжительность практики (количество недель) (только для практики,
% для остальных типов работ не используется)
\duration{17}

% Даты начала и окончания практики (только для практики, для остальных
% типов работ не используется)
\practStart{09.02.2026}
\practFinish{10.06.2026}

% Год выполнения отчета
\date{2026}

\maketitle
\tableofcontents

\section{Цель работы. Постановка задачи}

\textbf{Цель работы:} создание веб-приложения «Team Finder» --- платформы, на которой разработчики, дизайнеры и другие специалисты смогут находить единомышленников для совместной работы над pet-проектами. Зарегистрированные пользователи получат возможность размещать идеи проектов, собирать команду и откликаться на опубликованные предложения.

\textbf{Постановка задачи:} необходимо реализовать следующие функциональные возможности:

\begin{enumerate}
    \item Работа с пользователями. Создать кастомную модель пользователя с аутентификацией по email. Реализовать страницы регистрации, входа, выхода, просмотра и редактирования профиля, смены пароля.
    \item Пагинация. Подключить к проекту пагинацию и настроить вывод не более 12 проектов на главную страницу и на страницу списка пользователей.
    \item Аватарки пользователей. Добавить возможность загрузки аватарки. При регистрации должна автоматически генерироваться первоначальная аватарка с первой буквой имени.
    \item Создание проектов. У зарегистрированных пользователей должна быть возможность создавать новые проекты.
    \item Редактирование проектов. Автор проекта должен иметь возможность изменить его параметры.
    \item Система участия в проектах. Реализовать возможность присоединяться к проектам и отказываться от участия.
    \item Управление статусом и навыками проекта. Автор может завершить проект. Владелец проекта может добавлять и удалять навыки, необходимые проекту.
    \item Фильтрация проектов по навыкам. На странице списка проектов должна быть панель фильтрации по необходимым навыкам.
    \item Для выполнения задания необходимо использовать view-классы (CBV).
\end{enumerate}

\section{Реализация проекта}
\subsection{Работа с пользователями}

В проекте используется кастомная модель пользователя, унаследованная от \texttt{AbstractUser}, с аутентификацией по \texttt{email}.
Модель включает поля: \texttt{email} (уникальный, идентификатор), \texttt{first\_name}, \texttt{surname}, \texttt{avatar}, \texttt{phone}, \texttt{github\_url}, \texttt{about}.
Реализованы страницы регистрации, входа, выхода, просмотра и редактирования профиля, смены пароля.
Для регистрации используется собственная форма \texttt{UserRegisterForm}.

Модель \texttt{User} в \texttt{users/models.py}:

\begin{lstlisting}
class User(AbstractUser):
    email = models.EmailField(unique=True, verbose_name="Почта")
    first_name = models.CharField(
        max_length=constants.USER_FIRST_NAME_MAX_LENGTH, verbose_name="Имя")
    surname = models.CharField(
        max_length=constants.USER_SURNAME_MAX_LENGTH, verbose_name="Фамилия")
    avatar = models.ImageField(
        upload_to='avatars/', verbose_name="Аватар")
    phone = models.CharField(
        max_length=constants.PHONE_NUMBER_MAX_LENGTH,
        verbose_name='Номер телефона'
    )
    github_url = models.URLField(
        blank=True,
        default='',
        validators=[URLValidator()],
        verbose_name='Аккаунт на GitHub'
    )
    objects = UserManager()
    about = models.TextField(
        max_length=constants.USER_ABOUT_MAX_LENGTH,
        blank=True,
        default='',
        verbose_name='Описание профиля'
    )
    is_active = models.BooleanField(default=True)
    is_staff = models.BooleanField(default=False)
    USERNAME_FIELD = 'email'
    REQUIRED_FIELDS = ['first_name', 'surname']

    class Meta:
        verbose_name = 'Пользователь'
        verbose_name_plural = 'Пользователи'
        ordering = ['id']

    def __str__(self):
        return f"{self.first_name} {self.surname} ({self.email})"
\end{lstlisting}

Настройки в \texttt{users/urls.py}:

\begin{lstlisting}
from .views import (PasswordChangeView, UserDetailView, UserEditView,
                    UserListView, UserLoginView, UserLogoutView,
                    UserRegisterView)

app_name = 'users'

urlpatterns = [
    path('register/', UserRegisterView.as_view(), name='register'),
    path('login/', UserLoginView.as_view(), name='login'),
    path('logout/', UserLogoutView.as_view(), name='logout'),
    path('edit-profile/', UserEditView.as_view(), name='edit'),
    path('<int:user_id>/', UserDetailView.as_view(), name='user_detail'),
    path('change-password/', PasswordChangeView.as_view(), name='change-password'),
    path('list/', UserListView.as_view(), name='list'),
]
\end{lstlisting}

\begin{figure}[H]
    \centering
    \includegraphics[width=0.8\textwidth]{static/register.png}
    \caption{Форма для регистрации пользователя}
    \label{fig:registration}
\end{figure}

\begin{figure}[H]
    \centering
    \includegraphics[width=0.8\textwidth]{static/login.png}
    \caption{Форма для входа в систему}
    \label{fig:login}
\end{figure}

\begin{figure}[H]
    \centering
    
\includegraphics[width=0.8\textwidth]{static/user.png}
    \caption{Страница профиля пользователя}
    \label{fig:profile}
\end{figure}

\begin{figure}[H]
    \centering
    \includegraphics[width=0.8\textwidth]{static/user_edit.png}
    \caption{Страница редактирования профиля пользователя}
    \label{fig:profile_edit}
\end{figure}

\newpage
\subsection{Пагинация}

В проекте настроена пагинация на главной странице проектов и на странице списка пользователей.
Списки делятся на страницы по 12 элементов с помощью штатного механизма \texttt{ListView} и атрибута \texttt{paginate\_by}.
Константа \texttt{DEFAULT\_PAGINATION\_LIMIT} вынесена в отдельный файл \texttt{constants.py}.

View-класс \texttt{ProjectListView} \texttt{projects/views.py}:
 
\begin{lstlisting}
class ProjectListView(ListView):
    model = Project
    template_name = 'projects/project_list.html'
    context_object_name = 'projects'
    paginate_by = constants.DEFAULT_PAGINATION_LIMIT

    def get_queryset(self):
        queryset = Project.objects.select_related('owner').prefetch_related(
            'participants',
            'skills'
        )
        skill_name = self.request.GET.get('skill')
        if skill_name:
            queryset = queryset.filter(skills__name=skill_name.strip())
        return queryset

    def get_context_data(self, **kwargs):
        context = super().get_context_data(**kwargs)
        context['all_skills'] = Skill.objects.all()
        skill_name = self.request.GET.get('skill')
        if skill_name:
            context['active_skill'] = Skill.objects.filter(
                name=skill_name.strip()).first()
        else:
            context['active_skill'] = None
        return context
\end{lstlisting}

View-класс \texttt{UserListView} в \texttt{users/views.py}:

\begin{lstlisting}
class UserListView(ListView):
    model = User
    template_name = 'users/participants.html'
    context_object_name = 'participants'
    paginate_by = constants.DEFAULT_PAGINATION_LIMIT

    def get_context_data(self, **kwargs):
        context = super().get_context_data(**kwargs)
        context['query_prefix'] = ''
        return context
\end{lstlisting}

\begin{figure}[H]
    \centering
    \includegraphics[width=0.9\textwidth]{static/user_pagination.png}
    \caption{Пагинация списка проектов}
    \label{fig:project_pagination}
\end{figure}

На рисунках представлена страница списка проектов с пагинацией в размере 12 проектов на страницу.
Карточка проекта содержит заголовок, описание, информацию об авторе (ФИ и аватар), дату публикации и статус.
С помощью навигационной панели внизу страницы пользователь может перемещаться между страницами.

\newpage
\subsection{Аватарки пользователей}

В проекте реализована возможность загрузки аватарок пользователей.
В модель \texttt{User} добавлено поле \texttt{avatar} типа \texttt{ImageField}, настроена работа с файлами.
При создании пользователя генерируется первоначальная аватарка: первая буква имени на однотонном фоне со случайным цветом.
Файлы по умолчанию сохраняются в директории \texttt{media/avatars/}, указанной в настройках проекта.
Аватарки отображаются в шаблонах через тег \texttt{\{\{ user.avatar.url \}\}}.

Настройки в \texttt{users/models.py}:

\begin{lstlisting}
avatar = models.ImageField(upload_to='avatars/', verbose_name="Аватар")
\end{lstlisting}

Настройки в \texttt{team\_finder/settings.py}:

\begin{lstlisting}
MEDIA_URL = '/media/'
MEDIA_ROOT = BASE_DIR / 'media'
\end{lstlisting}

\begin{figure}[H]
    \centering
    
\includegraphics[width=0.8\textwidth]{static/image_post.png}
    \caption{Профиль пользователя с аватаркой}
    \label{fig:avatar}
\end{figure}

\newpage
\subsection{Добавление и редактирование проектов}

Пользователи имеют возможность создавать и редактировать проекты.
Классы \texttt{CreateView} и \texttt{UpdateView} обрабатывают создание и редактирование проектов.

При создании автоматически владельцем проекта устанавливается текущий пользователь, который также добавляется в участники.
При редактировании проверяется, что пользователь является владельцем проекта; при нарушении вызывается \texttt{PermissionDenied}.

Класс \texttt{ProjectCreateForm} в \texttt{projects/forms.py}:

\begin{lstlisting}
class ProjectCreateForm(forms.ModelForm):
    class Meta:
        model = Project
        fields = ('name', 'description', 'github_url', 'status')
        widgets = {
            'name': forms.TextInput(attrs={
                'placeholder': 'Название проекта'
            }),
            'description': forms.Textarea(attrs={
                'placeholder': 'Описание проекта',
                'rows': constants.PROJECT_FORM_DESCRIPTION_LENGTH
            }),
            'github_url': forms.URLInput(attrs={
                'placeholder': 'https://github.com/username/repo'
            }),
            'status': forms.Select(attrs={
                'class': 'form-select',
            }),
        }
\end{lstlisting}

View-класс \texttt{ProjectCreateView} в \texttt{projects/views.py}:

\begin{lstlisting}
class ProjectCreateView(LoginRequiredMixin, CreateView):
    model = Project
    form_class = ProjectCreateForm
    template_name = 'projects/create-project.html'
    success_url = reverse_lazy('projects:project_list')

    def form_valid(self, form):
        form.instance.owner = self.request.user
        response = super().form_valid(form)
        self.object.participants.add(self.request.user)
        return response

    def get_context_data(self, **kwargs):
        context = super().get_context_data(**kwargs)
        context['is_edit'] = False
        return context
\end{lstlisting}

View-класс \texttt{ProjectEditView} в \texttt{projects/views.py}:

\begin{lstlisting}
class ProjectEditView(LoginRequiredMixin, UpdateView):
    model = Project
    form_class = ProjectCreateForm
    template_name = 'projects/create-project.html'
    pk_url_kwarg = 'project_id'

    def get_success_url(self):
        return reverse_lazy('projects:project_detail',
                            kwargs={'project_id': self.object.pk})

    def dispatch(self, request, *args, **kwargs):
        obj = self.get_object()
        if obj.owner != request.user:
            raise PermissionDenied(
                'Вы не можете редактировать чужой проект!')
        return super().dispatch(request, *args, **kwargs)

    def get_context_data(self, **kwargs):
        context = super().get_context_data(**kwargs)
        context['is_edit'] = True
        return context
\end{lstlisting}

\begin{figure}[H]
    \centering
    \includegraphics[width=0.8\textwidth]{static/create_post.png}
    \caption{Форма создания проекта}
    \label{fig:create_project}
\end{figure}

\begin{figure}[H]
    \centering
    \includegraphics[width=0.8\textwidth]{static/edit_post.png}
    \caption{Форма редактирования проекта}
    \label{fig:edit_project}
\end{figure}

Используется Django ModelForm для автоматической генерации полей на основе модели \texttt{Project}.
В форме присутствует валидация поля \texttt{github\_url}: ссылка должна быть валидной и вести на GitHub.
Поле \texttt{status} представлено выпадающим списком с вариантами <<Открыт>> (open) и <<Закрыт>> (closed).

\newpage
\subsection{Участие в проектах}

В проекте реализована система участия в проектах.
Авторизованные пользователи могут присоединяться к чужим проектам или отказываться от участия.
Связь между пользователем и проектом реализована через ManyToMany-поле \texttt{participants} в модели \texttt{Project}.
Взаимодействие реализовано через POST-запросы с JSON-ответом.

Модель \texttt{Project} в \texttt{projects/models.py}:

\begin{lstlisting}
class Project(models.Model):
    class Status(models.TextChoices):
        OPEN = 'open', 'Открыт'
        CLOSED = 'closed', 'Закрыт'

    owner = models.ForeignKey(
        settings.AUTH_USER_MODEL,
        on_delete=models.CASCADE,
        related_name='owned_projects',
        verbose_name='Автор проекта'
    )
    name = models.CharField(
        max_length=constants.PROJECT_NAME_MAX_LENGTH,
        verbose_name='Название проекта')
    description = models.TextField(
        blank=True, verbose_name='Описание проекта')
    participants = models.ManyToManyField(
        settings.AUTH_USER_MODEL,
        blank=True,
        related_name='participated_projects',
        verbose_name='Участники')
    github_url = models.URLField(
        blank=True, verbose_name='Ссылка на GitHub')
    created_at = models.DateTimeField(
        auto_now_add=True, blank=True, verbose_name='Дата создания')
    status = models.CharField(
        max_length=constants.PROJECT_STATUS_MAX_LENGTH,
        choices=Status.choices,
        default=Status.OPEN,
        verbose_name='Статус')
    skills = models.ManyToManyField(
        Skill,
        blank=True,
        related_name='projects',
        verbose_name='Навыки проекта')

    class Meta:
        verbose_name = 'Проект'
        verbose_name_plural = 'Проекты'
        ordering = ['-created_at']

    def __str__(self):
        return self.name
\end{lstlisting}

View-класс \texttt{AddParticipantView} в \texttt{projects/views.py}:

\begin{lstlisting}
class AddParticipantView(LoginRequiredMixin, View):
    def post(self, request, project_id):
        project = get_object_or_404(Project, pk=project_id)

        if project.status != Project.Status.OPEN:
            return JsonResponse({
                'status': 'error',
                'message': 'В завершенный проект нельзя вступить'
            }, status=HTTPStatus.BAD_REQUEST)

        if project.owner == request.user:
            return JsonResponse({
                'status': 'error',
                'message': 'Вы являетесь владельцем своего проекта'
            }, status=HTTPStatus.BAD_REQUEST)

        if project.participants.filter(pk=request.user.pk).exists():
            project.participants.remove(request.user)
            response_data = {
                'status': 'ok',
                'participant': None
            }
        else:
            project.participants.add(request.user)
            response_data = {
                'status': 'ok',
                'participant': {
                    'id': request.user.id,
                    'name': f'{request.user.first_name} {request.user.surname}',
                    'avatar': (
                        request.user.avatar.url
                        if request.user.avatar else None
                    )
                }
            }
        return JsonResponse(response_data)
\end{lstlisting}

\begin{figure}[H]
    \centering
    \includegraphics[width=0.8\textwidth]{static/comments_create.png}
    \caption{Страница проекта с кнопкой участия и список участников проекта}
    \label{fig:participate}
\end{figure}

\newpage
\subsection{Закрытие проектов}

Реализование закрытие владельцем проекта, владелец может сменить его статус на <<closed>>.

View-класс \texttt{ProjectCompleteView} в \texttt{projects/views.py}:

\begin{lstlisting}
class ProjectCompleteView(LoginRequiredMixin, View):
    def post(self, request, project_id):
        project = get_object_or_404(Project, pk=project_id)

        if project.owner != request.user:
            return JsonResponse(
                {'status': 'error',
                 'message': 'Вы не можете завершить чужой проект!'},
                status=HTTPStatus.FORBIDDEN
            )

        if project.status != Project.Status.OPEN:
            return JsonResponse(
                {'status': 'error',
                 'message': 'Проект уже завершён'},
                status=HTTPStatus.CONFLICT
            )

        project.status = Project.Status.CLOSED
        project.save()
        return JsonResponse({'status': 'ok'})
\end{lstlisting}

\newpage
\subsection{Управление навыками}

Также реализовано управление навыками проекта: владелец может добавлять и удалять навыки без перезагрузки страницы.
При добавлении навыка работает автодополнение: GET-запрос к \texttt{projects/skills/?q=...} возвращает первые 10 подходящих навыков из базы.

View-классы в \texttt{skills/views.py}:

\begin{lstlisting}
class SkillAutocompleteView(View):
    def get(self, request):
        query = request.GET.get('q', '')
        skills = Skill.objects.filter(
            name__icontains=query
        )[:constants.SKILLS_PAGINATION_LIMIT]
        data = list(skills.values('id', 'name'))
        return JsonResponse(data, safe=False)


class AddSkillToProjectView(LoginRequiredMixin, View):
    def post(self, request, project_id):
        project = get_object_or_404(Project, id=project_id)
        if project.owner != request.user:
            return JsonResponse(
                {'error': 'У вас нет прав для редактирования этого проекта!'},
                status=HTTPStatus.FORBIDDEN
            )
        try:
            data = json.loads(request.body)
        except json.JSONDecodeError:
            return JsonResponse(
                {'error': 'Некорректный формат JSON'},
                status=HTTPStatus.BAD_REQUEST
            )
        skill_id = data.get('skill_id')
        name = data.get('name')
        if not skill_id and not name:
            return JsonResponse(
                {'error': 'Необходимо передать skill_id или name'},
                status=HTTPStatus.BAD_REQUEST
            )
        if skill_id:
            skill = get_object_or_404(Skill, id=skill_id)
        else:
            skill, _ = Skill.objects.get_or_create(name=name.strip())
        project.skills.add(skill)
        return JsonResponse({'id': skill.id, 'name': skill.name},
                            status=HTTPStatus.OK)


class RemoveSkillFromProjectView(LoginRequiredMixin, View):
    def post(self, request, project_id, skill_id):
        project = get_object_or_404(Project, pk=project_id)
        if project.owner != request.user:
            return JsonResponse(
                {'error': 'У вас нет прав для редактирования этого проекта!'},
                status=HTTPStatus.FORBIDDEN
            )
        skill = get_object_or_404(Skill, pk=skill_id)
        if project.skills.filter(id=skill.id).exists():
            project.skills.remove(skill)
            return JsonResponse({'status': 'ok'})
        return JsonResponse(
            {'status': 'error', 'message': 'Навыка нет в проекте'},
            status=HTTPStatus.BAD_REQUEST
        )
\end{lstlisting}

\newpage
\section{Заключение}

В ходе выполнения проекта было создано веб-приложение «Team Finder», представляющее собой платформу для поиска единомышленников для совместной работы над pet-проектами.

Выполнены все задачи с использованием современных возможностей фреймворка Django: кастомная модель пользователя с аутентификацией по email, пагинация, загрузка аватарок, создание и редактирование проектов, управление участниками и навыками проектов, фильтрация по навыкам, завершение проектов, использование CBV.

\newpage
\section*{СПИСОК ИСПОЛЬЗОВАННЫХ ИСТОЧНИКОВ}
\addcontentsline{toc}{section}{СПИСОК ИСПОЛЬЗОВАННЫХ ИСТОЧНИКОВ}

\begin{enumerate}
    \item Лутц Марк. Изучаем Python [Текст] / Марк Лутц --- Санкт-Петербург :  ООО <<Диалектика>>, 2019.
    \item Любанович Билл. Простой Python. Современный стиль программирования [Текст] / Билл Любанович --- Санкт-Петербург :  ООО Издательство <<Питер>>, 2016.
    \item Customizing authentication in Django [Электронный ресурс]. --- [Б. м. : б. и.]. --- URL: \url{https://docs.djangoproject.com/en/5.2/topics/auth/customizing/} (дата обращения: 01.06.2026). Загл. с экр. Яз. англ.

    \item Form and field validation [Электронный ресурс]. --- [Б. м. : б. и.]. --- URL: \url{https://docs.djangoproject.com/en/5.2/ref/forms/validation/} (дата обращения: 01.06.2026). Загл. с экр. Яз. англ.
    
    \item Paginator [Электронный ресурс]. --- [Б. м. : б. и.]. --- URL: \url{https://docs.djangoproject.com/en/5.2/ref/paginator/} (дата обращения: 01.06.2026). Загл. с экр. Яз. англ.

    \item Class-based views [Электронный ресурс]. --- [Б. м. : б. и.]. --- URL: \url{https://docs.djangoproject.com/en/5.2/topics/class-based-views/} (дата обращения: 01.06.2026). Загл. с экр. Яз. англ.

    \item QuerySet API reference [Электронный ресурс]. --- [Б. м. : б. и.]. --- URL: \url{https://docs.djangoproject.com/en/5.2/ref/models/querysets/} (дата обращения: 01.06.2026). Загл. с экр. Яз. англ.

\end{enumerate}

\end{document}
